\begin{table}
\centering
\caption{Kinematic properties of KURVS-CDFS galaxies.}
\label{tab2:kinematics}
\begin{tabular}{cccccc} 
\hline
\hline
KURVS ID & R$_{\rm H\alpha, max}$  & $\sigma_0$ & $v_{\rm rot}/\sigma_0$ & $t$  \\
 & kpc & km s$^{-1}$  & km s$^{-1}$ &  &  \\
  (1) & (2) & (3) & (4)  & (5) \\
\hline
1 & 10.7 & 96 $\pm$ 2 & 0.4 $\pm$ 0.1 & 1.05 $\pm$ 0.03 & \\
2 & 12.1 & 47 $\pm$ 1 & 1.3 $\pm$ 0.3 & 1.13 $\pm$ 0.12 & \\
3 & 11.7 & 51 $\pm$ 4 & 3.9 $\pm$ 0.1 & 0.92 $\pm$ 0.03 & \\
4 & 12.0 & 43 $\pm$ 2 & 0.2 $\pm$ 0.1 & 0.71 $\pm$ 0.05 & \\
5 & 12.1 & 44 $\pm$ 3 & 1.2 $\pm$ 0.2 & 1.19 $\pm$ 0.07 & \\
6 & 8.1 & 64 $\pm$ 13 & 1.4 $\pm$ 0.2 & 1.13 $\pm$ 0.10 & \\
7 & 11.4 & 61 $\pm$ 5 & 1.7 $\pm$ 0.2 & 0.98 $\pm$ 0.04 & \\
8 & 11.0 & 40 $\pm$ 3 & 1.8 $\pm$ 0.4 & 0.91 $\pm$ 0.03 & \\
9 & 10.2 & 48 $\pm$ 4 & 3.0 $\pm$ 0.3 & 1.04 $\pm$ 0.06 & \\
10$^{\star}$ & 15.2 & 61 $\pm$ 2 & 1.9 $\pm$ 0.1 & 0.99 $\pm$ 0.02 & \\
11 & 12.4 & 62 $\pm$ 2 & 4.4 $\pm$ 0.2 & 1.15 $\pm$ 0.04 & \\
12 & 11.2 & 127 $\pm$ 9 & 1.3 $\pm$ 0.3 & 0.67 $\pm$ 0.21 & \\
13 & 12.1 & 62 $\pm$ 4 & 1.5 $\pm$ 0.1 & 1.11 $\pm$ 0.05 & \\
14 & 9.8 & 102 $\pm$ 3 & 1.0 $\pm$ 0.1 & 1.28 $\pm$ 0.07 & \\
15 & 9.2 & 68 $\pm$ 4 & 1.7 $\pm$ 0.1 & 1.03 $\pm$ 0.03 & \\
16 & 10.9 & 69 $\pm$ 4 & 2.0 $\pm$ 0.1 & 1.02 $\pm$ 0.02 & \\
17 & 9.9 & 65 $\pm$ 3 & 1.7 $\pm$ 0.1 & 1.33 $\pm$ 0.12 & \\
18 & 12.1 & 61 $\pm$ 5 & 1.2 $\pm$ 0.1 & 1.11 $\pm$ 0.05 & \\
19 & 7.9 & 80 $\pm$ 4 & 1.2 $\pm$ 0.2 & 1.41 $\pm$ 0.38 & \\
20 & 10.4 & 51 $\pm$ 2 & 1.0 $\pm$ 0.3 & 1.18 $\pm$ 0.06 & \\
21 & 12.7 & 74 $\pm$ 2 & 1.8 $\pm$ 0.1 & 1.33 $\pm$ 0.08 & \\
22 & 13.6 & 155 $\pm$ 7 & 0.2 $\pm$ 0.1 & 0.65 $\pm$ 0.04 & \\
\hline
\hline
\end{tabular}
\\[2mm] 
\begin{flushleft}
{\bf Note.} 
(1) KURVS ID;
(2) Maximal extent of the observed rotation curve;
(3) Intrinsic velocity dispersion, corrected for instrumental broadening and beam smearing (see text for details);
(4) Rotational support;
(5) Outer rotation curve slope.
$^{\star}$ {While this galaxy displays the typical kinematics signatures of a rotating disc in KURVS observations, the {\it HST} imaging suggests that these are associated with the orbital motion of a merging pair (see Appendix \ref{A:KURVS data} for details). 
% is formally classified as a rotationally-supported disc as it displays a continuous velocity gradient along the kinematic major axis, has $v_{\rm rot}/\sigma_{0} \geqslant 1.5$ and the position of the steepest velocity gradient coincides with the peak of the velocity dispersion.However, the {\it HST} imaging suggests that these kinematics signatures are associated with the orbital motion of a merging pair. This is expected, given that that there is a high probability of misclassifying merging pairs as discs in seeing-limited observations \citep{Simons19, Rizzo22}.
We therefore exclude this galaxy from the dark matter fraction analysis presented in Sections \ref{Sect:Inner_RCs} and \ref{Sect:fDM}.}
 
\end{flushleft}
\end{table}
